% Extracted from p10-benchmark
\subsection{TinyLoop-Bench: A Benchmark Proposal}
\label{ch-bench}

\paragraph{Introduction}

Benchmarks shape fields. GLUE told us what ``understanding'' meant for
an era; perplexity leaderboards told us what fluency meant. For
recurrent language models there is no equivalent instrument. A
researcher who trains a looped transformer today reports the numbers
available---perplexity at the trained loop count---and the
loop-specific questions go unanswered: does quality improve
monotonically with loops? Where does it saturate? How does a 44M
recurrent model compare, per parameter, against a 44M fixed-depth
model?

A companion study in this series ran exactly one such evaluation: a
test-time loop sweep on a 44M OpenMythos model (43,974,722 parameters;
dim 512; 6 prelude + 2 coda layers; 4 loops; 8 experts; GQA; GPT-2
BPE 50,257; TinyStories, 120M tokens, loss 9.5242$\rightarrow$1.5263),
measuring perplexity at 1, 2, 4, 8, 12, and 16 loops. The result---a
sharp drop to 4.41 at the trained depth, then a flat plateau---was
obtained with ad-hoc scripts. TinyLoop-Bench exists to make that
evaluation standard, repeatable, and comparable across models.

\paragraph{Design Principles}

\begin{enumerate}
\item \textbf{Loops are first-class.} Every metric is reported as a
function of loop count, not just at the trained depth. A model card
without a loop-scaling curve is incomplete.
\item \textbf{Cheap to run.} Sub-1B models exist so that iteration is
fast; the suite must run on a single consumer GPU in under an hour.
No task may require more than 512 evaluation sequences.
\item \textbf{Proposed, versioned, frozen.} The suite ships with a
version number and frozen evaluation scripts; results cite the
version. This paper specifies v0.1.
\item \textbf{Honest about construct validity.} TinyStories-scale
data \citep{eldan2023tinystories} measures fluency and basic
coherence, not reasoning. The suite claims nothing beyond its
constructs.
\end{enumerate}

\paragraph{Task Categories}

\paragraph{C1: Fluency.} Standard language modelling on held-out
TinyStories sequences (256 sequences $\times$ 256 tokens, fixed
split, seed-pinned). Metric: perplexity at each loop count
$n \in \{1, 2, 4, 8, 12, 16\}$. This is the baseline everything else
builds on.

\paragraph{C2: Coherence over long generations.} The model generates
continuations of 256 tokens from 64 fixed prompts. Two proposed
metrics:
\begin{itemize}
\item \emph{Entity consistency}: named entities introduced in the
prompt are extracted (by a fixed, open NER pass---proposed
implementation: spaCy \texttt{en\_core\_web\_sm}, pinned version);
the score is the fraction still referenced consistently (no
contradictory attributes) in the continuation. Reported per loop
count.
\item \emph{Repetition rate}: fraction of generated 8-grams that
repeat a previous 8-gram within the continuation. Lower is better;
reported per loop count.
\end{itemize}
Both are heuristic proxies, documented as such; the suite reports
them alongside, never instead of, human-readable samples (5 fixed
prompts, outputs published verbatim in the report).

\paragraph{C3: Loop-scaling curves.} The signature category. For each
model, the suite records the full perplexity-vs-loops curve and
derives three summary statistics (definitions proposed here):
\begin{itemize}
\item \emph{Refinement gain} $G = \mathrm{PPL}(1)/\mathrm{PPL}(K_{\mathrm{train}})$:
how much iteration helps at all.
\item \emph{Saturation point} $n^*$: the smallest $n$ with
$\mathrm{PPL}(n) \leq 1.02 \times \min_m \mathrm{PPL}(m)$---where
gains effectively stop.
\item \emph{Extrapolation slope} $E = \mathrm{PPL}(2K_{\mathrm{train}})
- \min_m \mathrm{PPL}(m)$: positive means over-looping hurts, near
zero means it plateaus, negative means genuine extrapolation.
\end{itemize}
For the companion 44M model these would read (illustrative, from the
companion study's real numbers): $G = 29.38/4.41 \approx 6.7$,
$n^* = 4$, $E = 4.90 - 4.41 = +0.49$---refinement helps enormously,
saturates exactly at the trained depth, and extrapolation slightly
hurts. The suite's job is to make triples like $(6.7, 4, +0.49)$
the standard way recurrent models are compared.

\paragraph{C4: Parameter-efficiency frontiers.} Each submitted model is
plotted against fixed-depth baselines at matched parameter counts on
the C1--C3 metrics: quality per parameter. The suite ships reference
baselines (proposed: 44M and 130M fixed-depth transformers trained on
the same TinyStories split with the same token budget---to be trained
when the suite is built). A recurrent model ``wins'' its frontier if
it dominates the baseline curve at equal or fewer parameters.

\paragraph{Reference Evaluation Protocol}

\paragraph{Environment.} Single GPU, fp16 or bf16, batch size chosen
to fit 12\,GB VRAM; the harness auto-selects. All randomness seeded;
seeds published.

\paragraph{Procedure.} (1) Load checkpoint and config; the harness
reads the trained loop count $K_{\mathrm{train}}$ from config.
(2) Run C1 at $n \in \{1,2,4,8,12,16\}$ (capped at $4K_{\mathrm{train}}$
if larger). (3) Run C2 generations at $n \in \{1, K_{\mathrm{train}},
2K_{\mathrm{train}}\}$. (4) Compute C3 statistics from the C1 curve.
(5) Emit a JSON report plus a human-readable model card following the
fixed template (see below).

\paragraph{Reporting template.} Every submission reports: model
identifier, parameter count, architecture summary (dim, layers,
loops, experts, attention type, vocabulary), training tokens and
dataset, the full C1 table, the $(G, n^*, E)$ triple, C2 scores,
C4 frontier position, suite version, and the five verbatim
generations. No cherry-picking: the harness generates the card, not
the author.

\paragraph{Harness Design (Proposed)}

A single Python package, \texttt{tinyloop-bench}, with three parts:
\emph{tasks} (dataset loaders, generation loop, metric functions),
\emph{runners} (a model interface requiring only
\texttt{generate(prompt, n\_loops)} and \texttt{perplexity(text,
n\_loops)}), and \emph{reports} (JSON + Markdown card). The model
interface is deliberately minimal so that any recurrent
implementation---OpenMythos-style \citep{gomez2024openmythos} or
otherwise---can plug in. Reference datasets are content-addressed by
SHA-256; the exact TinyStories split and prompts ship with the
package.

\paragraph{Discussion}

The most likely objection is scope: TinyStories is toy data, and a
benchmark built on it cannot speak to reasoning or knowledge. We
accept this and bound the claims accordingly---but note that the
loop-scaling question is \emph{cleanest} at small scale, where full
sweeps are cheap and confounders are few. A suite that standardises
$(G, n^*, E)$ at 44M creates the vocabulary in which larger-scale
recurrent evaluation can later be discussed.

A second risk is Goodhart drift: once $(G, n^*, E)$ becomes a target,
architectures may optimise the triple rather than genuine refinement
(e.g.\ by making loop 1 artificially bad to inflate $G$). The
verbatim generations and the C4 frontier are partial mitigations;
ultimately, versioned frozen tasks plus community scrutiny are the
defence, as with every benchmark.

\paragraph{Limitations}

This paper specifies a benchmark; it does not build or run one. The
metric definitions (entity consistency, $G$, $n^*$, $E$) are
proposals awaiting validation---in particular, whether entity
consistency as defined correlates with human judgements of coherence
is untested. Reference baselines do not yet exist. The suite covers
only decoder-only recurrent transformers on English story data in
v0.1.

\paragraph{Conclusion}

Recurrent models need recurrent benchmarks. TinyLoop-Bench v0.1
specifies four task categories, concrete metrics including the
loop-scaling triple $(G, n^*, E)$, a frozen evaluation protocol, and
a harness design---everything needed to build it. The next step is
building it.
